\chapter{Conception et architecture}

\section{Vue d'ensemble}
L'architecture retenue se base sur une separation claire entre:
\begin{itemize}[leftmargin=1.2cm]
  \item un frontend React execute sur le poste utilisateur,
  \item un backend FastAPI execute sur le bastion,
  \item un cluster OpenShift cible pour l'hebergement final des VMs.
\end{itemize}

Cette separation repond a une contrainte forte: certaines VMs locales ne sont pas directement accessibles depuis le bastion.

\section{Architecture logicielle}
Le projet est structure autour des modules suivants:
\begin{itemize}[leftmargin=1.2cm]
  \item \texttt{src/api}: endpoints REST,
  \item \texttt{src/discovery}: decouverte des VMs,
  \item \texttt{src/analysis}: analyse de compatibilite,
  \item \texttt{src/conversion}: planification de conversion,
  \item \texttt{src/migration}: orchestration des migrations,
  \item \texttt{src/openshift}: integration avec les outils cluster,
  \item \texttt{src/ml}: recommandation basee IA,
  \item \texttt{frontend/frontend-app}: interface de pilotage.
\end{itemize}

\section{Choix techniques}
\begin{itemize}[leftmargin=1.2cm]
  \item Backend: FastAPI (Python)
  \item Frontend: React + Vite
  \item Outils de migration: \texttt{oc}, \texttt{virtctl}, \texttt{qemu-img}
  \item Cible: OpenShift Virtualization (KubeVirt)
\end{itemize}

\section{Diagrammes a inserer}
\textbf{A_COMPLETER:}
\begin{itemize}[leftmargin=1.2cm]
  \item diagramme d'architecture globale,
  \item diagramme de sequence du pipeline Analyze -> Plan -> Migrate,
  \item diagramme de classes (optionnel) pour les modules backend.
\end{itemize}
